\paragraph{В.1. Основы комбинаторики}

\subparagraph{В.1.1}
Сколько $k$-подстрок имеется у $n$-строки? (Одинаковые подстроки, начинающиеся в разных позициях строки, 
считаются разными.) Сколько всего подстрок имеется у строки длиной $n$?

\textbf{Ответ:}

В $n$-строке имеется $n-k+1$ различных $k$- подстрок.

Всего имеется $$\sum_{k=1}^{n} (n-k+1) = \sum_{j=1}^{n} j = \frac{n(n+1)}{2}$$

\subparagraph{В.1.2}
Булева функция (boolean function) с $n$ входами и $m$ выходами — это функция с областью определения
$\{\text{TRUE, FALSE}\}^n$ и областью значений $\{\text{TRUE, FALSE}\}^m$.
Сколько всего имеется различных функций с $n$ входами и одним выходом? С $n$ входами и $m$ выходами?

\textbf{Ответ:}

\textbf{С одним выходом ($m=1$):} 
    Всего существует $2^n$ различных входных комбинаций. Для каждой комбинации есть 2 варианта значения выхода. 
    Следовательно, общее количество функций: \[ 2^{2^n} \]

\textbf{С $m$ выходами:} 
    Всего существует $2^n$ входных комбинаций. Для каждой из них существует $2^m$ возможных вариантов выходных значений 
    (всевозможные битовые строки длины $m$). Таким образом, общее количество функций:
    \[ (2^m)^{2^n} = 2^{m \cdot 2^n} \]

\subparagraph{В.1.3}
Сколькими способами $n$ профессоров могут разместиться на конференции за круглым столом? 
Варианты, отличающиеся поворотом, считаются одинаковыми.

\textbf{Ответ:}

«Развернём» стол в линию, установив начальное и конечное положения.
Тогда на первое место может сесть любой из $n$ профессоров, на второе место — любой из оставшихся $n-1$ профессоров и так далее.
Таким образом, количество линейных перестановок равно:
\[
n! = n \cdot (n-1) \cdot (n-2) \cdots 1
\]
Однако для каждой такой рассадки существует $n$ поворотов стола, которые по условию считаются одинаковыми. 
Следовательно, общее количество уникальных способов размещения профессоров равно:
\[
\frac{n!}{n} = (n-1)!
\]

\subparagraph{В.1.4}
Сколькими способами можно выбрать из множества $\{1, 2, \dots, 99\}$ три различных числа так,
чтобы их сумма была четной?

\textbf{Ответ:}

В множестве $\{1, 2, \dots, 99\}$ содержится 49 четных чисел и 50 нечетных. Сумма трех чисел будет четной в двух случаях:
\begin{enumerate}
    \item Все три выбранных числа четные. Количество способов:
    \[ \binom{49}{3} = \frac{49 \cdot 48 \cdot 47}{6} = 18\,424 \]
    \item Одно число четное и два числа нечетные. Количество способов:
    \[ \binom{49}{1} \cdot \binom{50}{2} = 49 \cdot \frac{50 \cdot 49}{2} = 60\,025 \]
\end{enumerate}
Общее количество способов:
\[ 18\,424 + 60\,025 = 78\,449 \]

\subparagraph{В.1.5}
Докажите тождество
\[ \binom{n}{k} = \frac{n}{k} \binom{n-1}{k-1} \]
для $0 < k \le n$.

\textbf{Ответ:}

Раскроем формулу слева и справа:
\[ 
\frac{n}{k} \binom{n-1}{k-1} = \frac{n}{k} \cdot \frac{(n-1)!}{(k-1)!((n-1)-(k-1))!} = 
\]
\[ 
= \frac{n}{k} \cdot \frac{(n-1)!}{(k-1)!(n-k)!} = 
\]
\[ 
= \frac{n \cdot (n-1)!}{k \cdot (k-1)! \cdot (n-k)!} = \frac{n!}{k!(n-k)!} = \binom{n}{k}
\]
